\section{Newton systems: elimination, packing, and access}
\label{newt:section}

Cone resources do not determine the cost of a Newton equation.  The
projected variables, the auxiliary iterate, the chosen elimination, and
the requested output all matter.  We establish this point for the grouped
and column-packed ball formulations.  In particular, feasible auxiliary
allocations improve the off-center comparison from a quadratic to a
sharp linear dependence on residual eccentricity.

\subsection{The common marginal barrier and its Newton equation}

Let $C=\prod_{a=1}^b B_2^{s_a}$, let $n=\sum_a s_a$, and partition each
source vector $x_a$ into $h_a\geq1$ nonempty coordinate groups.  Pad a
group by public zero coordinates when necessary, and arrange the resulting
columns in real matrices $W_\ell\in\R^{p_\ell\times c_\ell}$.  Every
projected coordinate occurs once.  Let $E_a^\ell$ be the diagonal
projector onto the columns of block $\ell$ assigned to source $a$.
These projectors are orthogonal and sum to the identity in each block;
source $a$ has $h_a$ columns in total.  The packed formulation is
\begin{equation}\label{newt:packing}
 Z_\ell=\begin{pmatrix}S_\ell&W_\ell^T\\W_\ell&I_{p_\ell}\end{pmatrix}
 \succ0,\qquad
 \sum_\ell\tr(E_a^\ell S_\ell)=1,
 \qquad F=-\sum_\ell\log\det Z_\ell.
\end{equation}
All additional constraints and the objective in this subsection depend
only on $x$.  Write $D_\ell=S_\ell-W_\ell^TW_\ell$ and
$\rho_a=1-\norm{x_a}^2>0$.

\begin{proposition}[Packing-independent marginal and centered system]
\label{newt:center}
For fixed $x\in\intr C$, the unique minimum of $F$ on its auxiliary
fiber has diagonal residuals
\begin{equation}\label{newt:d0}
 (D_\ell^0)_{\gamma\gamma}=d_a:=\rho_a/h_a
       \quad(\gamma\hbox{ assigned to }a).
\end{equation}
The marginal barrier and its Hessian are
\begin{align}
 \Phi(x)&=-\sum_a h_a\log\rho_a+\sum_a h_a\log h_a,
       \label{newt:marginal}\\
 H_0(x)&=\bigoplus_a\left(
       \frac{2h_a}{\rho_a}I_{s_a}
       +\frac{4h_a}{\rho_a^2}x_ax_a^T\right).
       \label{newt:center-hessian}
\end{align}
Its self-concordant barrier parameter is exactly $\sum_a h_a$.
The grouped paraboloid formulation with the same groups has the same
marginal.  At their fiber centers, both formulations give the same
projected Newton matrix and right-hand side for every common affine
constraint $Ax=d$, linear objective, multiplier, and projected iterate.
\end{proposition}
\begin{proof}
The Schur determinant identity gives $\det Z_\ell=\det D_\ell$.
For fixed diagonal, Hadamard's inequality uniquely maximizes the
determinant at a diagonal positive matrix.  The allocation equations become
$\sum_{\gamma\in\Gamma_a}d_\gamma=\rho_a$; AM--GM uniquely maximizes
their product at $d_\gamma=\rho_a/h_a$.  The same AM--GM calculation
applies to the grouped slacks $t_\gamma-\norm{x_\gamma}^2$.
Differentiation gives \eqref{newt:center-hessian}.  Each summand is $h_a$
times the standard unit-ball barrier, and its squared dual gradient norm is
$2h_a\norm{x_a}^2/(1+\norm{x_a}^2)$.  This tends to $h_a$ at the
boundary.  Additivity and simultaneous boundary approach prove exactness
of the parameter.

For completeness, the equality extends to the uncondensed residual
equations.  For a direction $(A_\ell,U_\ell)=(\delta S_\ell,\delta W_\ell)$
put $B_\ell=A_\ell-W_\ell^TU_\ell-U_\ell^TW_\ell$.  Direct
differentiation along this affine direction gives
\begin{equation}\label{newt:lifted-hessian}
 D^2F[(A,U)]^2=\sum_\ell\left[
 \tr(D_\ell^{-1}B_\ell D_\ell^{-1}B_\ell)
       +2\tr(D_\ell^{-1}U_\ell^TU_\ell)\right].
\end{equation}
At diagonal $D^0$, off-diagonal entries of $B$ are independent positive
quadratic modes.  They have zero gradient, objective, and constraint
coefficients.  The diagonal modes $r_\gamma=B_{\gamma\gamma}$ obey
\begin{equation}\label{newt:star-quadratic}
 \sum_\gamma\frac{r_\gamma^2}{d_\gamma^2}
       +2\sum_\gamma\frac{\norm{u_\gamma}^2}{d_\gamma},
 \qquad
 \sum_{\gamma\in\Gamma_a}
       (r_\gamma+2x_\gamma^Tu_\gamma)=0.
\end{equation}
These are precisely the grouped residual equations, including their
linear terms $-\sum_\gamma r_\gamma/d_\gamma$ and the common objective.
Additional rows $Au=0$ coincide.  Eliminating $r$ and the allocation
multipliers therefore gives the same matrix and right-hand side.
\end{proof}

With no additional affine coupling, the graph of
\eqref{newt:star-quadratic}, with one allocation multiplier per source,
is a disjoint union of stars.  Solving these stars, or applying the
diagonal-plus-rank-one inverse in \eqref{newt:center-hessian}, costs
$O(n+\sum_a h_a)=O(n)$ arithmetic.  The graph of a fully materialized
generic Hessian is instead a union of $s_a$-cliques.  Both descriptions
are correct; they measure different representations of the same operator.
Adding constraint hubs for $A$ gives a graph $G$ on projected variables,
source hubs, and constraint hubs.  Attaching the $r_\gamma$ leaves gives
width $\max\{1,\operatorname{tw}(G)\}$; this graph, rather than the PSD
packing layout, controls the corresponding sparse solve.

\begin{corollary}[Matched projected-oracle equivalence]
\label{newt:oracle}
Supply both formulations with the same complete canonical oracles for
their common projected data and reduced Newton equation, at the same
precision.  Require the same projected direction, or its normalized
quantum state when the direction is nonzero.  Then their minimum
classical and quantum query complexities agree exactly.  This equivalence
also holds along any adaptive procedure that uses only these projected
records and recenters the auxiliary fibers before each Newton call.
\end{corollary}
\begin{proof}
Proposition~\ref{newt:center} identifies both the operator and the right-hand
side, in the same coordinates.  Forward each query to the identical oracle;
retain all other operations, measurements, and stopping decisions.  The
simulation works in both directions.  Induction on calls gives the
adaptive statement.  A common charged compiler from common raw data
likewise gives identical preparation work.
\end{proof}

Here ``same oracle'' includes its action on unused input and ancilla
subspaces.  Equality of an encoded matrix block alone would permit a
unitary completion to reveal extra information.  Thus the corollary is
an exact interface statement, not an oracle construction at no cost.
It preserves normalization, spectrum, right-hand-side norm, solution
overlap, and output accuracy.  In residual coordinates the centered PSD
KKT system is the grouped system direct-summed with a positive diagonal
off-diagonal block and a zero right-hand side; public zero padding also
identifies those solution states.  The original $\delta S$ output is
different: reconstructing it requires $W^TU+U^TW$.  Its cost is not covered
by the projected equivalence.  If all $h_a=h$, changing $h$ multiplies
the unconstrained marginal Hessian by $h$.  It changes its metric scale
but not its condition number or normalized inverse state for a fixed
right-hand side.  This illustrates why fewer cone factors alone do not
establish a state-generation advantage.

\subsection{A sharp comparison away from the fiber center}

\begin{theorem}[Exact quotient and allocation-aware transfer]
\label{newt:linear}
At a strictly feasible point of \eqref{newt:packing}, define
\begin{align}
 z_a(U)&=\sum_\ell\tr(E_a^\ell W_\ell^TU_\ell),
 &M_{ab}(D)&=\sum_\ell\tr(E_a^\ell D_\ell E_b^\ell D_\ell).
       \label{newt:z-m}
\end{align}
The quadratic form obtained by eliminating all auxiliary Hessian
directions is
\begin{equation}\label{newt:quotient}
 Q_D(U)=2\sum_\ell\tr(D_\ell^{-1}U_\ell^TU_\ell)
                 +4z(U)^TM(D)^{-1}z(U).
\end{equation}
Let $H_D$ denote this form on the projected coordinates.  If
$\mu D_\ell^0\preceq D_\ell\preceq L D_\ell^0$ for all $\ell$,
where $0<\mu\leq1\leq L$, then
\begin{equation}\label{newt:linear-transfer}
 L^{-1}H_0\preceq H_D\preceq\mu^{-1}H_0,
 \qquad
 \kappa(H_0^{-1/2}H_DH_0^{-1/2})\leq L/\mu.
\end{equation}
The dependence and the condition factor are sharp.  The comparison
survives restriction to any common projected tangent space.  For any
full-row-rank matrix $A$, the dual normal matrices satisfy
\begin{equation}\label{newt:dual-transfer}
 \mu AH_0^{-1}A^T\preceq AH_D^{-1}A^T
                         \preceq L AH_0^{-1}A^T.
\end{equation}
\end{theorem}
\begin{proof}
Equip symmetric residual directions with
$\ip{B}{C}_D=\sum_\ell\tr(D_\ell^{-1}B_\ell D_\ell^{-1}C_\ell)$.
The allocation functional $B\mapsto\sum_\ell\tr(E_a^\ell B_\ell)$
has Riesz representative $(D_\ell E_a^\ell D_\ell)_\ell$.  The Gram
matrix of these representatives is $M(D)$.  It is positive definite:
for $E_\ell(\theta)=\sum_a\theta_aE_a^\ell$,
\[
 \theta^TM(D)\theta=\sum_\ell
   \norm{D_\ell^{1/2}E_\ell(\theta)D_\ell^{1/2}}_F^2,
\]
and its vanishing forces every $\theta_a=0$, since each source has a
column.  The tangent allocation constraint is
$\sum_\ell\tr(E_a^\ell B_\ell)=-2z_a(U)$.
Its minimum squared $D$-norm is $4z^TM^{-1}z$, proving
\eqref{newt:quotient} from \eqref{newt:lifted-hessian}.

The sharper comparison uses feasibility, not just tensor-product order.
Although $E_\ell(\theta)$ may be indefinite,
$E_\ell(\theta)D_\ell E_\ell(\theta)\succeq0$.  Positive trace pairing
therefore gives
\[
 \mu\sum_\ell\tr(D_\ell^0 E_\ell D_\ell E_\ell)
 \leq\theta^TM(D)\theta
 \leq L\sum_\ell\tr(D_\ell^0 E_\ell D_\ell E_\ell).
\]
The diagonal matrices $D_\ell^0$ and $E_\ell$ commute.  Exact source
allocations, $\sum_\ell\tr(E_a^\ell D_\ell)=\rho_a$, yield
\begin{equation}\label{newt:allocation-identity}
 \sum_\ell\tr(D_\ell^0 E_\ell D_\ell E_\ell)
 =\sum_a\theta_a^2d_a\rho_a
 =\theta^TM_0\theta,
 \qquad M_0=\operatorname{diag}(\rho_a^2/h_a).
\end{equation}
Hence $\mu M_0\preceq M(D)\preceq LM_0$.
Both terms in \eqref{newt:quotient} now obey the first-power bounds
$L^{-1}$ and $\mu^{-1}$.  At the same $W$, their vector $z(U)$ is
unchanged, so these inequalities prove \eqref{newt:linear-transfer}.
Restriction preserves order; inversion followed by congruence proves
\eqref{newt:dual-transfer}.  Sharpness is established in
Example~\ref{newt:obstruction} below.
\end{proof}

It is essential to distinguish this Hessian elimination from first
minimizing the barrier value over a fiber.  The latter always replaces
$D$ by $D^0$.  It is also essential to retain the eliminated gradient.
For minimizing $F+\eta\ip{c}{x}$ at a feasible packed point, the reduced
Newton gradient is the linear functional
\begin{equation}\label{newt:reduced-gradient}
 g_D[U]=2\rho^TM(D)^{-1}z(U)+\eta\ip{c}{U}.
\end{equation}
Indeed the Newton quadratic before auxiliary elimination is
\[
 \tfrac12\norm{B}_D^2-\sum_\ell\tr(D_\ell^{-1}B_\ell)
  +\sum_\ell\tr(D_\ell^{-1}U_\ell^TU_\ell)+\eta\ip{c}{U}.
\]
Completing the square replaces $B$ by $B-D$ and its allocation by
$-2z(U)-\rho$.  Its minimum is
$\tfrac12(2z+\rho)^TM^{-1}(2z+\rho)$, plus the displayed $U$ terms
and a constant.  This proves \eqref{newt:reduced-gradient}.  At $D^0$
it is $\sum_a2h_a\ip{x_a}{U_a}/\rho_a+\eta\ip{c}{U}$,
the marginal gradient.  Away from the center, a spectral comparison of
matrices alone does not identify Newton directions or their normalized
states.  Homogeneous additional tangent constraints are common to both
systems; nonzero affine infeasibility has additional right-hand-side terms.

\begin{example}[Sharpness and fixed-projection obstruction]
\label{newt:obstruction}
Pack two sources as two columns of one block, take $W=0$, and put
\[
 D=\begin{pmatrix}1&\xi\\\xi&1\end{pmatrix},\qquad
 D^0=I_2,\qquad |\xi|<1.
\]
The source equations hold exactly.  Formula~\eqref{newt:quotient} reduces
to
\[
 Q_D(U)=\frac{2}{1-\xi^2}
  (\norm{u_1}^2+\norm{u_2}^2-2\xi u_1^Tu_2).
\]
The relative eigenvalues are $1/(1+\xi)$ and $1/(1-\xi)$.
Thus $\mu=1-|\xi|$, $L=1+|\xi|$, and the relative condition is
exactly $L/\mu$.  For objective columns $(v,0)$, $\norm{v}=1$, the
projected Newton directions are $-\eta(v,\xi v)/2$ and
$-\eta(v,0)/2$ at $D$ and $D^0$.  Their normalized states have trace
distance $|\xi|/\sqrt{1+\xi^2}$.  Replacing $\xi$ by $-\xi$
leaves all projected input records unchanged but makes the two states
have trace distance $2|\xi|/(1+\xi^2)\to1$.
Consequently projected data alone cannot uniformly recover unrestricted
off-center Newton states.  This is a one-iterate obstruction; it asserts
no lower bound on the number of outer iterations.
\end{example}

\begin{corollary}[Approximate allocations]
\label{newt:inexact}
Keep $W$ and its exact center $D^0$ fixed, but allow positive residuals
$\widehat D$ with allocations
\[
 (1-\zeta)\rho_a\leq
 \sum_\ell\tr(E_a^\ell\widehat D_\ell)
 \leq(1+\zeta)\rho_a,\qquad 0\leq\zeta<1.
\]
If $\mu D^0\preceq\widehat D\preceq LD^0$, with
$0<\mu\leq1\leq L$, their homogeneous-tangent Hessian quotient obeys
\begin{equation}\label{newt:approx-transfer}
 \frac{H_0}{L(1+\zeta)}\preceq H_{\widehat D}
       \preceq\frac{H_0}{\mu(1-\zeta)}.
\end{equation}
In particular relative residual error
$\norm{(D_\ell^0)^{-1/2}(\widehat D_\ell-D_\ell^0)
(D_\ell^0)^{-1/2}}_2\leq\delta<1$ gives condition at most
$((1+\delta)/(1-\delta))^2$ without exact allocations, and at most
$(1+\delta)/(1-\delta)$ when the allocations are exact.
\end{corollary}
\begin{proof}
The middle trace in \eqref{newt:allocation-identity} lies between
$(1-\zeta)\theta^TM_0\theta$ and
$(1+\zeta)\theta^TM_0\theta$.  Repeat the proof and use the weaker
of the resulting bounds for the two nonnegative terms of $Q$.
Spectral error $\delta$ implies allocation error at most $\delta$ by
positive trace pairing with each projector.  Set $\mu=1-\delta$,
$L=1+\delta$, and $\zeta=\delta$ or $0$, respectively.
\end{proof}

The linear improvement in Theorem~\ref{newt:linear} is specific to the
feasible allocation structure.  Its proof uses elementary positive trace
inequalities; it is not a new general matrix-order principle.  Neither
version supplies free access to $D^{-1}$, $M(D)^{-1}$, or $H_0^{-1/2}$.
Those constructions and any encoding normalization remain separate costs.

\subsection{What exact recentering requires}

The center has an inexpensive implicit representation.  Store $(W,d)$
with the semantics $S_\ell=W_\ell^TW_\ell+\operatorname{Diag}(d_\ell)$.
Computing all squared source norms and all labels
$d_\gamma=\rho_a/h_a$ uses $O(n)$ arithmetic and storage, independently
of packing, because $\sum_a h_a\leq n$.  The source equations hold
exactly by $\norm{x_a}^2+h_ad_a=1$.  For $B$-bit common-denominator
dyadic coordinates the squared norms and labels are exact rationals of
$O(B+\log n)$ bits.  Fast integer arithmetic gives the conservative
bound $\widetilde O(n\mathsf M(B+\log n))$, where $\mathsf M(t)$
is the cost of multiplying $t$-bit integers.  This represents the exact
center of the stored interior rational iterate.  It does not require
subtracting rounded Gram matrices.  Unrelated rational denominators,
inverse evaluation, and requested numerical accuracy have their own bit
costs.

A reversible value oracle for $\rho_a$ similarly supplies $d_\gamma$
with one scale query and reversible division by the public $h_a$.
Together with blockwise direction-state access, it permits the common
rank-one compiler for \eqref{newt:center-hessian}; radii and certified
normalization bounds must be supplied or computed.  In particular
normalized state access to $x_a$ alone does not provide this scale.
The same complete state-preparation unitary can prepare the direction of
$r_0v$ and $r_1v$, while their $1-r_j^2$ differ by arbitrarily large
factors.  No algorithm with only that unitary can return both scales to
a fixed relative accuracy.

\begin{proposition}[Cost of missing radial metadata]
\label{newt:scale}
For coordinate-bit access to a source of dimension $s\geq2$, producing a
classical center scale to relative error $\delta<1/3$ can require
$\Omega(\sqrt s)$ quantum or $\Omega(s)$ randomized queries.
For $b$ independent sources, returning all scales with joint success at
least $2/3$ requires $\Omega(b\sqrt s)$ and $\Omega(bs)$,
respectively.  These orders are attained on the promise used in the proof.
For $k\mid s$, $1\leq k\leq s/2$, the one-source promise can instead
have smaller slack $\rho=k/s$, with tight costs
$\Theta(\rho^{-1/2})$ quantum and $\Theta(\rho^{-1})$ randomized.
\end{proposition}
\begin{proof}
Take $x_i=z_i/\sqrt s$ with $z_1=0$ public and all other bits one,
except for at most one additional hidden zero.  The slack is $1/s$ or
$2/s$.  Their relative-$\delta$ intervals are disjoint precisely when
$\delta<1/3$, so recovering the scale solves promised unstructured
search.  The standard search bounds apply \citep{BBBV1997,BHMT2002}.

For clarity, the quantum direct sum has a short adversary proof.  Index
one source's promise by the no-mark input and its $s-1$ single-mark
inputs.  The star adjacency matrix $\Gamma$ has norm $\sqrt{s-1}$,
and filtering by any queried bit leaves norm one.  For $b$ sources use
$\sum_{a=1}^b I^{\otimes(a-1)}\otimes\Gamma\otimes
I^{\otimes(b-a)}$.  Its norm is $b\sqrt{s-1}$; filtering a query in
one source retains only that summand.  Its entries vanish between inputs
having the same output vector.  The adversary lower bound for function
evaluation \citep{LMRSS2011} gives the claim.  For randomized algorithms,
give each source an independent fair no-mark/single-uniform-mark promise.
Joint success implies constant average success for each output, and the
usual search indistinguishability argument requires a constant fraction
of the $s-1$ expected queries per such source; summing proves $\Omega(bs)$.
On this promise, exact amplitude amplification for the known single-mark
case followed by verification distinguishes zero from one mark with
$O(\sqrt s)$ queries and no error.  Applied sourcewise it attains joint
success without an amplification logarithm.  Scanning attains the
randomized order.

Finally partition coordinates into blocks of size $k$.  One public block
is zero; all candidate blocks are one except possibly one hidden zero
block.  The two slacks are $k/s$ and $2k/s$.  A representative coordinate
is an oracle for its block mark, reducing to search on $s/k-1$ labels.
This number is within a constant factor of $s/k$ in the stated range.
\end{proof}

The joint-output bound also applies to a standalone compiled value
interface whose output can be read on all labels with the same joint
correctness guarantee and with no further coordinate-input queries.
It does not describe one online coherent evaluation: the promised search
can be controlled by the source label at $O(\sqrt s)$ input queries per
evaluation.  Those queries remain charged on every invocation.  The
following fresh update family supplies a temporal bound under an explicit
output contract.

A \emph{compile-and-commit maintainer} returns, before the next epoch,
either a classical record of all source scales or a complete reusable
unitary for their encoded numerical values.  The latter must allow every
basis-label value to be copied out, using the unitary and its inverse,
with no further raw-update queries and with the entire backing workspace
restored exactly, including correlations with other retained registers.
More precisely, each committed object has fixed value strings $v_a$;
compute--copy--uncompute on label $a$ acts as
\[
 |a,z\rangle|\psi\rangle\longmapsto
 |a,z\mathbin\oplus v_a\rangle|\psi\rangle
\]
on its backing state and any retained reference.  The strings encode
numbers $\widehat d_a$.  Correctness means that, with joint probability at
least $2/3$, all committed values at all epochs obey
$|\widehat d_a-d_a|\leq\delta d_a$.  A state consumed or disturbed by
reading its values does not meet this contract.  Rational encodings allow
exact values even for $\delta=0$; fixed-point encodings require enough
bits for the requested positive error.

\begin{theorem}[Fresh-update scale maintenance]
\label{newt:dynamic-scale}
Let $b,T\geq1$, $s\geq2$, and $0\leq\delta<1/3$.  For the packing
\eqref{newt:packing} with $b$ sources of dimension $s$ and public column
counts $h_a$, there is a family of $T$ fresh coordinate-bit update batches,
each zero or one-sparse per source, on which compile-and-commit maintenance
has tight worst-case raw-query costs
\begin{equation}\label{newt:dynamic-commit}
 Q=\Theta(Tb\sqrt s),\qquad R=\Theta(Tbs)
\end{equation}
for quantum and randomized algorithms, respectively.  Each epoch starts
from the same public interior base and is reset after commitment.
For $k\mid s$, $1\leq k\leq s/2$, a replicated-block version has
$k$-sparse updates and smaller slack $\rho_*=k/s$, with costs
\begin{equation}\label{newt:dynamic-margin}
 Q=\Theta(Tb\rho_*^{-1/2}),\qquad R=\Theta(Tb\rho_*^{-1}).
\end{equation}
The lower bounds allow all epoch oracles from the start, superposed queries
over epoch, source, and coordinate, and arbitrary persistent workspace.
The upper bounds have zero error on these promises.
\end{theorem}
\begin{proof}
Use $\bar x=(0,1,\ldots,1)/\sqrt s$.  Independently for each pair
$(t,a)$, the update is either $u_{t,a}=0$ or
$u_{t,a}=-e_j/\sqrt s$ at one hidden $j\in\{2,\ldots,s\}$.
The raw oracle returns the corresponding mark bit at a queried coordinate.
After the update,
\[
 \rho_{t,a}=1-\norm{\bar x+u_{t,a}}^2\in\{1/s,2/s\},
 \qquad d_{t,a}=\rho_{t,a}/h_a.
\]
Both cases are strictly interior, and Proposition~\ref{newt:center}
gives these exact PSD fiber-center labels.  Their relative-error intervals
are disjoint because $1+\delta<2(1-\delta)$.
Undo each batch after commitment; inverse-update access uses the same
bits and reveals no new information.  A public reset only strengthens
the model.

Read every committed label before continuing the maintainer, restoring its
workspace in the quantum case.  Thresholding recovers $m=Tb$ independent
presence bits with joint success $2/3$ and no extra raw queries.  Apply the
star adversary from Proposition~\ref{newt:scale}, with tensor factors now
indexed by $(t,a)$.  Its Kronecker sum has norm $m\sqrt{s-1}$ and every
coordinate filter has norm one.  It is entrywise nonnegative and vanishes
between inputs with the same output vector, so the finite-output spectral
adversary bound applies \citep[Theorems~3--4]{ACGT2010}.  The proof already
uses simultaneous oracle access and allows cross-epoch query scheduling.

For the randomized bound, independently choose no mark with probability
$1/2$ and a uniform single mark otherwise.  Fix the randomness and all
other pairs, and let $L$ count distinct candidates queried when this pair
has no mark.  The no-mark and random-mark transcripts coincide unless
the mark lies in these $L$ locations.  Averaging over the fixed data, the
total variation distance is at most $\mathbb E L/(s-1)$.
Marginal success at least $2/3$
forces $\mathbb E L\geq(s-1)/3$.  Since the no-mark case has probability
$1/2$, this pair contributes at least $(s-1)/6$ expected queries under
the product distribution.  Summing over the $m$ pairs proves the
worst-case randomized bound.

Exact amplitude amplification designed for one marked coordinate,
followed by verification, decides zero versus one mark in $O(\sqrt s)$
queries \citep[Theorems~4 and~16]{BHMT2002}.  Verification always rejects
in the zero-mark case.  Search each pair and store its exact rational
scale; no joint-success amplification logarithm is needed.  Scanning
gives the classical upper bound.

For the margin version, partition coordinates into $s/k$ blocks of size
$k$, with one public zero block in the base.  An update zeros either no
candidate block or one whole candidate block.  The slacks are $k/s$ and
$2k/s$.  A coordinate query reveals only its block's mark bit, and one
representative coordinate per block suffices.  Replace $s-1$ throughout
by $s/k-1=\Theta(s/k)$, which is at least one in the stated range.
\end{proof}

\begin{corollary}[Universal invocation versus compilation tradeoff]
\label{newt:scale-service}
On the same promises, suppose a service must support any adaptive client
making at most $r_t$ value invocations at epoch $t$, including coherent
label inputs in the quantum case.  Invocations implement the reusable
value interface above, but may query the raw input; all preprocessing and
invocation queries are charged.  Require joint correctness with probability
at least $2/3$ over the entire interaction.  For prescribed nonnegative
integer budgets $r_t$, the tight worst-case costs are
\begin{equation}\label{newt:scale-service-cost}
 Q=\Theta\left(\sqrt s\sum_{t=1}^T\min\{r_t,b\}\right),\qquad
 R=\Theta\left(s\sum_{t=1}^T\min\{r_t,b\}\right).
\end{equation}
In the margin version replace $\sqrt s,s$ by
$\rho_*^{-1/2},\rho_*^{-1}$, respectively.
\end{corollary}
\begin{proof}
Choose the permitted client that reads $\min\{r_t,b\}$ distinct basis
labels at each epoch.  Fix all other update pairs to zero.  Its transcript
recovers $\sum_t\min\{r_t,b\}$ independent search bits, so the same
adversary and product-distribution proofs give the lower bounds.
For the upper bound, if $r_t<b$, run the exact promised search controlled
by the source label on each invocation.  Coherently verify, copy the
deterministic presence bit, and uncompute the search workspace; this is a
clean value unitary at $O(\sqrt s)$ queries per call.  If $r_t\geq b$,
search all sources once and keep a classical table for query-free reversible
lookup.  Classical scans give the analogous strategies.  Apply the same
choice to candidate-block representatives for the margin version.
\end{proof}

This universal-service contract is stronger than a particular downstream
algorithm's structured calls.  One online coherent evaluation has no $b$
factor; a service supporting at least $b$ arbitrary calls reaches the
compilation cost.  Neither result concerns arbitrary one-shot state output
or finite-precision gate complexity.

For an explicit sparse list, the exact identity
\begin{equation}\label{newt:scale-update}
 \rho_a(x+\alpha u)=\rho_a(x)-2\alpha x_a^Tu_a
                               -\alpha^2\norm{u_a}^2
\end{equation}
instead maintains scales in $O(\sum_a|\operatorname{supp}u_a|)$
arithmetic, given access to the current coordinates and a list of affected
sources.  In particular one-sparse lists cost $O(b)$ per epoch.
The lower bound charges discovery of hidden support, not sparsity itself.
Nor does a snapshot lower bound imply a factor $T$ for one fixed oracle:
if every round requests the same promised OR bit, compute and cache it
once in $\Theta(\sqrt s)$ quantum or $\Theta(s)$ randomized queries.
All $T$ subsequent requests then use no raw queries.

\begin{proposition}[Fixed central-path scale caching]
\label{newt:fixed-scale}
For the uncoupled product of balls, minimize
$\Phi(x)+\eta\sum_a c_a^Tx_a$ with fixed $c_a$ and $\eta\geq0$.
Put $q_a=\norm{c_a}$ and $\tau_a=\eta q_a/h_a$.  The unique center is
\begin{equation}\label{newt:fixed-path-scale}
 \rho_a(\eta)=\frac{2}{1+\sqrt{1+\tau_a^2}},\qquad
 x_a(\eta)=-\frac{\eta\rho_a(\eta)}{2h_a}c_a,\qquad
 d_a(\eta)=\frac{\rho_a(\eta)}{h_a}.
\end{equation}
For $s_a\geq2$, its source Hessian has radial-to-tangential eigenvalue ratio
$2/\rho_a(\eta)-1=\sqrt{1+\tau_a^2}$.
\end{proposition}
\begin{proof}
Stationarity is $2h_ax_a/\rho_a+\eta c_a=0$.  Substitution in
$\rho_a=1-\norm{x_a}^2$ gives
$\rho_a+\tau_a^2\rho_a^2/4=1$; its positive solution is
\eqref{newt:fixed-path-scale}, also when $q_a=0$ or $\eta=0$.
It is interior, and strict convexity of $\Phi$ gives uniqueness.
Equation~\eqref{newt:center-hessian} gives eigenvalues $2h_a/\rho_a$
and $2h_a/\rho_a+4h_a(1-\rho_a)/\rho_a^2$, proving the ratio
(equal to one at $x_a=0$).
\end{proof}

Thus supplied reversible access to the fixed norms $q_a$ evaluates a
requested scale with $O(1)$ norm queries and charged scalar arithmetic;
uncomputing norm scratch may require an inverse query.  Explicit objective
data permits all norms to be acquired once in $O(n)$ real arithmetic,
including square roots, and reused for every $\eta$.
Requested precision and reversible arithmetic still have their own costs.
Additional affine constraints or inexact neighborhood iterates need not
satisfy this stationarity equation.  The fresh batches in
Theorem~\ref{newt:dynamic-scale} have not been embedded in the Newton
trajectory of one fixed conic program: such an embedding and a contract
recovering its independent outputs would be needed to multiply a query
lower bound by an iteration count.  The closed-form uncoupled path witnesses
reuse of fixed data, and supplies no such embedding.

Explicit dense auxiliaries have a different output cost:
$\Omega(\sum_\ell c_\ell^2)$ scalar writes and ordinary Gram formation
cost $O(\sum_\ell p_\ell c_\ell^2)$.  An implicit ambient encoding can
avoid these writes: put $K_\ell=(W_\ell\ I_{p_\ell})$ and write
$Z_\ell=K_\ell^TK_\ell+\operatorname{Diag}(d_\ell,0)$.
Products of a supplied $\alpha_K$-normalized encoding of $K_\ell$ and
its adjoint, followed by a positive linear combination with an
$\alpha_d$-normalized diagonal encoding, give normalization
$\alpha_K^2+\alpha_d$.  This costs a constant number of supplied encoding
calls.  Constructing these encodings, their norms, and an inverse remains
charged; no statement about a flattened Gram-matrix output state follows.

\subsection{Sparse augmented Lorentz systems}

Low-rank augmentation of cone Hessians is classical, including in
ECOS and nonsymmetric-cone implementations
\citep{DomahidiChuBoyd2013,ChenGoulart2025}.  We use it to make the
classical comparator explicit for the formulations above.  A graph below
means the fixed structural graph of the augmented scalar matrix; special
numerical zeros may delete edges.  Its treewidth is the minimum, over
tree decompositions, of the largest bag size minus one.

For a Lorentz block $x=(t,z)$, let $q=t^2-\norm z^2>0$ and
$J=\operatorname{diag}(1,-I)$.  Its standard Hessian is
\begin{equation}\label{newt:one-hub}
 H_x=-\frac{2J}{q}+\frac{4(Jx)(Jx)^T}{q^2}
       =D_x+w_xw_x^T.
\end{equation}
Thus the equality-form Newton matrix and its one-hub augmentation are
\begin{equation}\label{newt:augmented}
 K_0=\begin{pmatrix}H&A^T\\A&0\end{pmatrix},\qquad
 \widetilde K_0=
 \begin{pmatrix}D&W&A^T\\W^T&-I&0\\A&0&0\end{pmatrix}.
\end{equation}
Here $W$ has one embedded $w_x$ column per cone.  Eliminating the negative
identity gives $K_0$, and the augmented right-hand side is $(g_x,0,g_y)$.
There is no assumption of a positive diagonal in $D$.

\begin{proposition}[Exact field-operation comparison]
\label{newt:treewidth}
Suppose $A\in\R^{p\times M}$ has full row rank, there are $k$ Lorentz
blocks, and a compact width-$\tau$ tree decomposition of the structural
graph of $\widetilde K_0$ is supplied.  Then a Newton equation can be
solved in $O((M+p+k)(\tau+1)^2)$ exact field operations, after
$O(M+\operatorname{nnz}A)$ coefficient formation.  Compactness here
means $O(M+p+k)$ bags; processing a larger supplied decomposition is an
additional input cost.
\end{proposition}
\begin{proof}
Positive definiteness of $H$ and full row rank of $A$ imply nonsingularity
of both matrices in \eqref{newt:augmented}.  Replace each vertex of each
bag by its row and column copies.  The resulting decomposition of the
matrix's row--column bipartite graph has width at most $2\tau+1$:
every off-diagonal entry has both endpoints in some original bag, each
diagonal entry is covered, and connectedness of each copy's bags is
preserved.  Apply the exact low-treewidth linear-system algorithm of
\citet[Corollary~3]{FurerHoppenTrevisan2025}.  Its bound is
$O((m+n)(w+1)^2)$ for a supplied compact width-$w$ bipartite
decomposition.  It allows nonzero pivots chosen during elimination and
algebraic cancellations, so the zero or negative diagonal entries here
cause no difficulty.  Formula~\eqref{newt:one-hub} uses only field
operations and gives the formation cost.
\end{proof}

A positive-diagonal version is useful when a reusable symmetric
factorization is wanted.  If $r=\norm z>0$ and $e=z/r$, put
\begin{equation}\label{newt:two-hub}
 D=\frac2q I,\quad
 a=\frac{\sqrt{2r(t+r)}}q\binom1{-e},\quad
 v=\frac{\sqrt{2r(t-r)}}q\binom1e.
 \qquad H_x=D+aa^T-vv^T.
\end{equation}
At $r=0$ take $a=v=0$.  The identity follows by restricting $H_x-D$
to the span of $(1,0)$ and $(0,e)$, where its matrix is
$4q^{-2}\left(\begin{smallmatrix}r^2&-tr\\-tr&r^2\end{smallmatrix}\right)$.
Moreover $v^TD^{-1}v=2r/(t+r)<1$, so $D-vv^T\succ0$.
With embedded update and downdate columns denoted by $U,V$, respectively,
the regularized augmentation is
\begin{equation}\label{newt:quasidefinite}
 \begin{pmatrix}
 D&V&U&A^T\\V^T&I&0&0\\U^T&0&-I&0\\A&0&0&-\varpi I
 \end{pmatrix},\qquad \varpi>0.
\end{equation}
Grouping $(x,V)$ as positive variables and $(U,y)$ as negative variables
shows symmetric quasidefiniteness: the first diagonal block has Schur
complement $D-VV^T\succ0$, and the second is negative definite.
Every symmetric permutation therefore admits an exact pivot-free
$LDL^T$ factorization \citep{Vanderbei1995}.  A supplied width-$\tau$
elimination order gives $O((M+p+k)(\tau+1)^2)$ real arithmetic,
$O((M+p+k)(\tau+1))$ storage, and width-linear subsequent solves.
Square roots in \eqref{newt:two-hub} distinguish this branch from the
field-operation result.

This is not an unconditional numerical stability statement.
Writing $K_\varpi=K_0-\varpi\operatorname{diag}(0,I)$, the resolvent
identity gives, for $d_0=K_0^{-1}g\ne0$ and
$\varpi\norm{K_0^{-1}}_2<1$,
\begin{equation}\label{newt:regularization-error}
 \frac{\norm{K_\varpi^{-1}g-d_0}}{\norm{d_0}}
 \leq\frac{\varpi\norm{K_0^{-1}}_2}
              {1-\varpi\norm{K_0^{-1}}_2}.
\end{equation}
Working precision, small pivots, growth, and residual certification still
need bounds.  Quasidefiniteness rules out exact zero-pivot breakdown, not
arbitrarily small pivots.  Also, two hubs cannot generically be reduced to
one while retaining a positive diagonal base: when $m\geq3$ and every
coordinate of $Jx$ is nonzero, off-diagonal consistency in
$H=D_++\epsilon aa^T$ forces $\epsilon=1$ and
$a_i^2=4(Jx)_i^2/q^2$.  The axial entry of $D_+$ is then $-2/q$.
This only proves minimality within that diagonal-plus-one-signed-rank-one
class.

The two-hub graph can be bounded from coarse incidence.  Make one vertex
for each cone and each equality row, joining them when the row uses a
coordinate of that cone.  If this graph has width $w$ and each scalar
column of $A$ has at most one nonzero, the expanded graph has width at
most $\max\{2w+1,3\}$.  Replace every cone vertex in a coarse bag by its
two hubs.  For each used coordinate attach a bag consisting of that
coordinate, its two hubs, and its single row, to a bag covering the
corresponding coarse edge.  For unused coordinates omit the row.  These
bags cover all edges and preserve the connected-bag condition.  Thus a
coarse forest gives width at most three regardless of cone dimension.

\subsection{Exact graph counts for the ball constructions}

For $N$ projected coordinates grouped into $k$ nonempty groups, set
$q_j=s_j-\norm{w_j}^2$ and impose $\sum_j s_j=1$.
The reduced barrier has Hessian blocks
\begin{equation}\label{newt:group-hessian}
 \operatorname{diag}(0,2I/q_j)
      +q_j^{-2}(1,-2w_j)(1,-2w_j)^T.
\end{equation}
Their one-hub graph has the paths
$\lambda-s_j-h_j$ and leaves $h_j-w_{ji}$.
It is a tree with $N+2k+1$ vertices and $N+2k$ edges.
The zero $s_j$ diagonal prevents an unqualified scalar leaf-pivot claim;
Proposition~\ref{newt:treewidth} instead supplies the exact solve.
Restricting \eqref{newt:two-hub} by
$(s,w)\mapsto((1+s)/2,(1-s)/2,w)$ gives positive base
$\operatorname{diag}(1/q,2I/q)$ and a positive downdated base by
congruence.  The two-hub reduced graph has $N+3k+1$ vertices,
$2N+3k$ edges, and width two.  Bags containing the two hubs and one
coordinate, together with bags $\{\lambda,s_j\}$, prove the upper
bound; each nonempty group supplies a cycle for the lower bound.

For comparison, retain all raw Lorentz coordinates $(u_j,v_j,w_j)$
and equations $u_j+v_j=1$, $\sum_j(u_j-v_j)=1$.
The exact structural counts are
\begin{center}
\begin{tabular}{lrrc}
\toprule
representation & vertices & edges & treewidth\\
\midrule
raw grouped, one hub & $N+4k+1$ & $N+6k$ & $2$\\
raw grouped, two hubs & $N+5k+1$ & $2N+8k$ & $3$\\
reduced grouped, one hub & $N+2k+1$ & $N+2k$ & $1$\\
reduced grouped, two hubs & $N+3k+1$ & $2N+3k$ & $2$\\
\bottomrule
\end{tabular}
\end{center}
To verify the raw widths, the one-hub graph has $K_{2,3}$ on axial
coordinates $u_j,v_j$ and the hub, local row, and global row.  Bags
$\{u_j,v_j,a\}$ for these three vertices $a$, plus hub--coordinate
leaf bags, give width two.  The two-hub graph has a subdivision of $K_4$:
its branch vertices are the two hubs and $u_j,v_j$; a row connects the
two axial vertices and one $w$ coordinate connects the two hubs.
A width-three decomposition uses the four branch vertices as a core,
the two axial vertices with the local and global rows as another bag,
and bags of the two hubs with each $w$ coordinate.  Different groups
join through the global row.  Counting incidences gives the edge totals.

The direct ball formulation $(t,x)\in Q_{N+1}$, $t=1$, has a one-hub
tree with $N+3$ vertices; eliminating $t$ gives a star with $N+1$
vertices.  A binary norm tree with $L=N-1$ internal nodes has $3L$
raw cone coordinates, $L$ hubs, $L-1$ link rows, and one root row.
It is connected with $5L$ vertices and $3L+2(L-1)+1=5L-1$ edges,
hence a tree.  Its rows are independent because each link or root
equality uses its own cone-coordinate slots.  These counts also explain
why the smooth coordinatewise construction ($k=N$) has raw width two
but reduced width one.  Smooth contact selections do not impose an
intrinsic width-two cost.

All these isolated-ball Newton equations have linear exact arithmetic
cost.  The ambient and restricted barrier parameters and the movement
bounds proved earlier remain different.  Thus replacing a large cone by
many small cones can increase the standard short-step certificate and
the bounded-Dikin movement cost without improving this linear solve.
For products of cone epigraphs, fixing their axial coordinates to one
reduces to products of balls; the earlier additive curvature bound forces
at least $k(s-1)$ copies of $Q_3$ to represent $Q_{s+1}^k$, even under
joint reformulation.  Their ambient product barrier has parameter at
least $2k(s-1)$ by its orthant section.  Separate binary trees attain the
count.  This compares the generic ambient-barrier certificate with the
$2k$ certificate for the direct product, not all algorithms' iteration
counts.

\subsection{The full-output comparison and its boundary}

The preceding exact solves give a precise replacement principle.  Suppose
an outer algorithm requests a classical direction or updated iterate in
every round, its theorem accepts any direction satisfying a stated Newton
residual contract, and the current coefficients and right-hand side are
classically evaluable at the cost charged to its quantum access model.
If fixed compact decompositions of widths $\tau_t$ are supplied for its
expanded systems, replacing the quantum Newton module by
Proposition~\ref{newt:treewidth} costs
\begin{equation}\label{newt:replacement}
 \sum_t O((M_t+p_t+k_t)(\tau_t+1)^2)
\end{equation}
exact field operations, plus coefficient formation and any decomposition
input cost.  The outer guarantee is preserved because the exact solution
meets its residual contract.  This need not reproduce the quantum
algorithm's random trajectory; the comparison requires the outer
guarantee and the structural width bound for every admitted trajectory.
Dense materialization already requires $\Omega(M_t+p_t)$ scalar writes
per round.  When $p_t+k_t=O(M_t)$ and $\tau_t=M_t^{o(1)}$, the
Newton module therefore gives no polynomial dimensional saving under
this matched full-output, exact-arithmetic contract.  The centered packed
family and the grouped or norm-tree ball family have $O(n)$ or $O(N)$
replacements.  Dense external affine coupling can invalidate their width
bounds.  A finite-precision replacement additionally needs a stable
implementation and certification at the prescribed tolerance.

There is a complementary classical route for sparse well-conditioned
systems.  For a consistent full-column-rank system $Kd=r$ and zero
initial guess, conjugate gradients on its normal equations, implemented
by matrix--vector products as CGLS or LSQR \citep{PaigeSaunders1982},
has the exact-arithmetic residual estimate
\[
 \norm{Kd_j-r}_2\leq
 2\left(\frac{\kappa_2(K)-1}{\kappa_2(K)+1}\right)^j\norm r_2.
\]
This follows from the usual CG energy-norm bound, because that norm is
$\norm{K(d_j-d_*)}_2$ and
$\sqrt{\kappa_2(K^TK)}=\kappa_2(K)$.
If one pair of products by $K,K^T$ costs $\mathcal M$, including current
entry evaluation, a relative residual $0<\epsilon_{\rm lin}<1$ costs
$O(\mathcal M\kappa_2(K)\log(2/\epsilon_{\rm lin}))$.
The case $\kappa_2(K)=1$ terminates after one step.  No normal matrix
need be materialized.  Sparse bounded-size conic blocks fit this
comparison only when their actual augmented entries are cheaply
available.  Quantum-only encodings and tall implicit Hessian formation
are different access models; in particular this argument does not
replace the row-sampling work in \citet{ApersGribling2026}.

The classical algorithms, signed-rank augmentations, and search bounds
used here are established tools.  The formulation-specific conclusions
are the packing-independent centered equation, its sharp feasible
off-center transfer, and the exact access and output conditions under
which the same classical Newton work applies.  Published quantum SDP
and portfolio IPM analyses explicitly account for conditioning, data
access, and classical readout \citep{AugustinoEtAl2023,DalzellEtAl2023};
the present comparison is confined to the displayed structured families.
In a block-diagonal SDP embedding the matrix order is the sum of block
orders, not their maximum.  Substituting a block cap for that sum in a
generic solver theorem is therefore unjustified.  State, sample, scalar,
or coherent-iterate output can have different costs, as the following
query results demonstrate.
